% Slug: trap-hoist
% Source: chapters/ch01/traps-effects.tex:231 (tikzpicture env start line)
% Original label: \label{fig:trap-hoist}
% Original caption (LaTeX source):
%   Hoisting a risky operation above its guard adds behaviors the source lacked.
% Original caption (plain text):
%   Hoisting a risky operation above its guard adds behaviors the source lacked.
% Required TikZ setup: \usepackage{tikz} + \usetikzlibrary{arrows.meta, positioning, shapes.geometric, fit} (loaded by support/compiler-pseudocode.sty).
% Shared styles used in this picture: `basic block', `cfg edge'.
%   Defined in support/compiler-pseudocode.sty: `basic block', `cfg edge', `back edge', `block label'.
% Usage: % Slug: trap-hoist
% Source: chapters/ch01/traps-effects.tex:231 (tikzpicture env start line)
% Original label: \label{fig:trap-hoist}
% Original caption (LaTeX source):
%   Hoisting a risky operation above its guard adds behaviors the source lacked.
% Original caption (plain text):
%   Hoisting a risky operation above its guard adds behaviors the source lacked.
% Required TikZ setup: \usepackage{tikz} + \usetikzlibrary{arrows.meta, positioning, shapes.geometric, fit} (loaded by support/compiler-pseudocode.sty).
% Shared styles used in this picture: `basic block', `cfg edge'.
%   Defined in support/compiler-pseudocode.sty: `basic block', `cfg edge', `back edge', `block label'.
% Usage: % Slug: trap-hoist
% Source: chapters/ch01/traps-effects.tex:231 (tikzpicture env start line)
% Original label: \label{fig:trap-hoist}
% Original caption (LaTeX source):
%   Hoisting a risky operation above its guard adds behaviors the source lacked.
% Original caption (plain text):
%   Hoisting a risky operation above its guard adds behaviors the source lacked.
% Required TikZ setup: \usepackage{tikz} + \usetikzlibrary{arrows.meta, positioning, shapes.geometric, fit} (loaded by support/compiler-pseudocode.sty).
% Shared styles used in this picture: `basic block', `cfg edge'.
%   Defined in support/compiler-pseudocode.sty: `basic block', `cfg edge', `back edge', `block label'.
% Usage: % Slug: trap-hoist
% Source: chapters/ch01/traps-effects.tex:231 (tikzpicture env start line)
% Original label: \label{fig:trap-hoist}
% Original caption (LaTeX source):
%   Hoisting a risky operation above its guard adds behaviors the source lacked.
% Original caption (plain text):
%   Hoisting a risky operation above its guard adds behaviors the source lacked.
% Required TikZ setup: \usepackage{tikz} + \usetikzlibrary{arrows.meta, positioning, shapes.geometric, fit} (loaded by support/compiler-pseudocode.sty).
% Shared styles used in this picture: `basic block', `cfg edge'.
%   Defined in support/compiler-pseudocode.sty: `basic block', `cfg edge', `back edge', `block label'.
% Usage: \input{graphs-code/trap-hoist.tikz.tex} inside a figure environment
%   (wrap with \begin{figure} ... \caption{...} \label{fig:trap-hoist} ... \end{figure}).
% Note: tikzpicture body below is verbatim from the source; do not reformat.
\begin{tikzpicture}[node distance=14mm and 12mm]
    \node[basic block] (test) {test $c$};
    \node[basic block,below left=of test] (yes) {risky $op$};
    \node[basic block,below right=of test] (no) {skip};
    \node[basic block,below right=of yes] (ret) {return $r$};
    \draw[cfg edge] (test) -- node[above left,font=\small] {true} (yes);
    \draw[cfg edge] (test) -- node[above right,font=\small] {false} (no);
    \draw[cfg edge] (yes) -- (ret);
    \draw[cfg edge] (no) -- (ret);
    \node[basic block,right=38mm of test] (op2) {risky $op$};
    \node[basic block,below=of op2] (test2) {test $c$};
    \node[basic block,below=of test2] (ret2) {return $r$};
    \draw[cfg edge] (op2) -- (test2);
    \draw[cfg edge] (test2) -- (ret2);
    \node[font=\sffamily\small] at (1.2,1.0) {Before};
    \node[font=\sffamily\small] at (6.6,1.0) {After};
  \end{tikzpicture}
 inside a figure environment
%   (wrap with \begin{figure} ... \caption{...} \label{fig:trap-hoist} ... \end{figure}).
% Note: tikzpicture body below is verbatim from the source; do not reformat.
\begin{tikzpicture}[node distance=14mm and 12mm]
    \node[basic block] (test) {test $c$};
    \node[basic block,below left=of test] (yes) {risky $op$};
    \node[basic block,below right=of test] (no) {skip};
    \node[basic block,below right=of yes] (ret) {return $r$};
    \draw[cfg edge] (test) -- node[above left,font=\small] {true} (yes);
    \draw[cfg edge] (test) -- node[above right,font=\small] {false} (no);
    \draw[cfg edge] (yes) -- (ret);
    \draw[cfg edge] (no) -- (ret);
    \node[basic block,right=38mm of test] (op2) {risky $op$};
    \node[basic block,below=of op2] (test2) {test $c$};
    \node[basic block,below=of test2] (ret2) {return $r$};
    \draw[cfg edge] (op2) -- (test2);
    \draw[cfg edge] (test2) -- (ret2);
    \node[font=\sffamily\small] at (1.2,1.0) {Before};
    \node[font=\sffamily\small] at (6.6,1.0) {After};
  \end{tikzpicture}
 inside a figure environment
%   (wrap with \begin{figure} ... \caption{...} \label{fig:trap-hoist} ... \end{figure}).
% Note: tikzpicture body below is verbatim from the source; do not reformat.
\begin{tikzpicture}[node distance=14mm and 12mm]
    \node[basic block] (test) {test $c$};
    \node[basic block,below left=of test] (yes) {risky $op$};
    \node[basic block,below right=of test] (no) {skip};
    \node[basic block,below right=of yes] (ret) {return $r$};
    \draw[cfg edge] (test) -- node[above left,font=\small] {true} (yes);
    \draw[cfg edge] (test) -- node[above right,font=\small] {false} (no);
    \draw[cfg edge] (yes) -- (ret);
    \draw[cfg edge] (no) -- (ret);
    \node[basic block,right=38mm of test] (op2) {risky $op$};
    \node[basic block,below=of op2] (test2) {test $c$};
    \node[basic block,below=of test2] (ret2) {return $r$};
    \draw[cfg edge] (op2) -- (test2);
    \draw[cfg edge] (test2) -- (ret2);
    \node[font=\sffamily\small] at (1.2,1.0) {Before};
    \node[font=\sffamily\small] at (6.6,1.0) {After};
  \end{tikzpicture}
 inside a figure environment
%   (wrap with \begin{figure} ... \caption{...} \label{fig:trap-hoist} ... \end{figure}).
% Note: tikzpicture body below is verbatim from the source; do not reformat.
\begin{tikzpicture}[node distance=14mm and 12mm]
    \node[basic block] (test) {test $c$};
    \node[basic block,below left=of test] (yes) {risky $op$};
    \node[basic block,below right=of test] (no) {skip};
    \node[basic block,below right=of yes] (ret) {return $r$};
    \draw[cfg edge] (test) -- node[above left,font=\small] {true} (yes);
    \draw[cfg edge] (test) -- node[above right,font=\small] {false} (no);
    \draw[cfg edge] (yes) -- (ret);
    \draw[cfg edge] (no) -- (ret);
    \node[basic block,right=38mm of test] (op2) {risky $op$};
    \node[basic block,below=of op2] (test2) {test $c$};
    \node[basic block,below=of test2] (ret2) {return $r$};
    \draw[cfg edge] (op2) -- (test2);
    \draw[cfg edge] (test2) -- (ret2);
    \node[font=\sffamily\small] at (1.2,1.0) {Before};
    \node[font=\sffamily\small] at (6.6,1.0) {After};
  \end{tikzpicture}
